\documentclass[11pt,a4paper]{article}
% preamble.tex — shared setup for human-verification documents
% Compile each document with:  tectonic <File>.tex

\usepackage[T1]{fontenc}
\usepackage[utf8]{inputenc}
\usepackage{lmodern}
\usepackage{amssymb,amsmath,amsthm}
\usepackage{mathtools}
\usepackage[dvipsnames]{xcolor}
\usepackage{listings}
\usepackage{hyperref}
\usepackage[margin=2.5cm]{geometry}
\usepackage{parskip}

% ── Colours ──────────────────────────────────────────────────────────
\definecolor{leanBg}{HTML}{F7F7F0}
\definecolor{haskellBg}{HTML}{F0F4F7}
\definecolor{keyword}{HTML}{7B3294}
\definecolor{comment}{HTML}{5A8C32}
\definecolor{noteFg}{HTML}{B35900}

% ── Lean 4 listings language ─────────────────────────────────────────
% NOTE: Lean listing blocks use ASCII approximations of unicode symbols
% (e.g. -> for arrow, o. for compose).  The literate key substitutes
% these with proper math-mode renderings in the PDF.
\lstdefinelanguage{Lean4}{
  morekeywords={def,theorem,structure,where,namespace,end,import,fun,
                let,have,show,by,exact,intro,cases,rfl,subst,obtain,
                funext,simp,rw,apply,variable,Type,Prop,instance,
                section,open,class,extends,deriving,sorry,noncomputable},
  sensitive=true,
  morecomment=[l]{--},
  morecomment=[n]{/-}{-/},
  morestring=[b]",
}

% ── Lean listing style ───────────────────────────────────────────────
\lstdefinestyle{lean}{
  language=Lean4,
  backgroundcolor=\color{leanBg},
  basicstyle=\small\ttfamily,
  keywordstyle=\color{keyword}\bfseries,
  commentstyle=\color{comment}\itshape,
  breaklines=true,
  frame=single,
  framerule=0.4pt,
  xleftmargin=4pt,
  xrightmargin=4pt,
  aboveskip=6pt,
  belowskip=4pt,
  literate=
    {->}{{$\to$}}2
    {=>}{{$\Rightarrow$}}2
    {<|}{{$\langle$}}1
    {|>}{{$\rangle$}}1
    {==>}{{$\Longrightarrow$}}3
    { o. }{{$\;\circ\;$}}3
    {forall}{{$\forall$}}1
    {:=}{{$\,:=\,$}}2
    {!=}{{$\neq$}}2
    {<=}{{$\leq$}}2
    {>=}{{$\geq$}}2,
}

% ── Haskell listing style ───────────────────────────────────────────
\lstdefinestyle{haskell}{
  language=Haskell,
  backgroundcolor=\color{haskellBg},
  basicstyle=\small\ttfamily,
  keywordstyle=\color{keyword}\bfseries,
  commentstyle=\color{comment}\itshape,
  breaklines=true,
  frame=single,
  framerule=0.4pt,
  xleftmargin=4pt,
  xrightmargin=4pt,
  aboveskip=6pt,
  belowskip=4pt,
  literate=
    {->}{{$\to$}}2
    {=>}{{$\Rightarrow$}}2,
}

% ── Theorem environments ────────────────────────────────────────────
\theoremstyle{definition}
\newtheorem{defn}{Definition}[section]
\newtheorem{thm}[defn]{Theorem}
\newtheorem{prop}[defn]{Proposition}

\theoremstyle{remark}
\newtheorem*{fidelitynote}{Fidelity note}

% ── Convenience macros ───────────────────────────────────────────────
\newcommand{\lean}[1]{\texttt{#1}}      % inline Lean identifier
\newcommand{\hask}[1]{\texttt{#1}}      % inline Haskell identifier
\newcommand{\Ext}{\mathrm{Ext}}
\newcommand{\Pos}{\mathrm{Pos}}
\newcommand{\Shape}{\mathrm{Shape}}


\title{Human Verification: \texttt{Containers.Network}\\[4pt]
  \large Shipping Networks as Directed Containers}
\author{Machine-generated from \texttt{lean/Network.lean} for review by Neil Ghani}
\date{July 2026}

\begin{document}
\maketitle

\begin{abstract}
This document presents every definition and theorem \emph{statement} from the
Lean~4 module \lean{Containers.Network} in three parallel forms: the original
Lean, a readable Haskell rendering, and precise mathematical English.  Proofs
are omitted --- they are machine-checked by Lean.  The purpose is to let a
human verifier confirm that the formalisation captures the intended
mathematics.

The file verifies Construction~\texttt{con:network} (\S4 of the paper): every
directed graph $G$ induces a directed container $\mathcal{N}(G)$ whose shapes
are nodes, whose positions at a node $v$ are finite directed paths out of $v$,
and whose directed-container laws follow from path-concatenation identities.
The file ends with the 11-facility Sawasdee Cellars wine network as the
paper's running example.
\end{abstract}

\tableofcontents
\bigskip

%===================================================================
\section{Overview}
%===================================================================

A \emph{digraph} $G = (V, E, \mathrm{src}, \mathrm{tgt})$ gives rise to a
directed container $\mathcal{N}(G)$ --- the \emph{network container} ---
following the recipe:
\begin{itemize}
  \item $\mathbf{Shape} = V$ (nodes);
  \item $\mathbf{Pos}(v) = \{(w, \rho) \mid \rho \text{ is a path } v \rightsquigarrow w\}$
    (routes out of $v$, together with their endpoint);
  \item $\mathbf{root}(v) = (v, \varepsilon)$ (the empty path);
  \item $\mathbf{sub}(v, (w, \rho)) = w$ (the endpoint of the route);
  \item $\mathbf{shift}(v, (w, \rho), (u, \sigma)) = (u, \rho \mathbin{+\!\!+} \sigma)$
    (route concatenation).
\end{itemize}
As a category (via \lean{DContCat.cat}) this is the free category on $G$.

%===================================================================
\section{Directed graphs}
%===================================================================

% --- Digraph -------------------------------------------------------
\subsection{Structure: \lean{Digraph}}

\begin{defn}[Directed graph]
A \textbf{directed (multi-)graph} $G$ consists of a type $V$ of vertices, a
type $E$ of edges, and two maps $\mathrm{src}, \mathrm{tgt} : E \to V$ giving
the source and target of each edge.  Parallel edges and loops are allowed; no
finiteness is assumed.
\end{defn}

\paragraph{Lean.}
\begin{lstlisting}[style=lean]
structure Digraph where
  V   : Type
  E   : Type
  src : E -> V
  tgt : E -> V
\end{lstlisting}

\paragraph{Haskell.}
\begin{lstlisting}[style=haskell]
data Digraph v e = Digraph
  { src :: e -> v
  , tgt :: e -> v
  }
\end{lstlisting}


% --- Path ----------------------------------------------------------
\subsection{Inductive type: \lean{Digraph.Path}}

\begin{defn}[Directed path]
A \textbf{directed path} $\rho : \mathrm{Path}_G(v, w)$ is defined
inductively:
\begin{itemize}
  \item $\mathrm{nil}(v) : \mathrm{Path}_G(v, v)$ --- the empty path at $v$;
  \item $\mathrm{cons}(e, \rho) : \mathrm{Path}_G(\mathrm{src}(e), w)$ when
    $\rho : \mathrm{Path}_G(\mathrm{tgt}(e), w)$ --- prepend an edge.
\end{itemize}
Paths are cons-at-the-front: the first constructor argument is an outgoing
edge from the start node, followed by a path from its target.
\end{defn}

\paragraph{Lean.}
\begin{lstlisting}[style=lean]
inductive Path (G : Digraph) : G.V -> G.V -> Type
  | nil  (v : G.V)                            : Path G v v
  | cons (e : G.E) {w : G.V}
         (p : Path G (G.tgt e) w)             : Path G (G.src e) w
\end{lstlisting}

\paragraph{Haskell.}
\begin{lstlisting}[style=haskell]
-- Approximation: a path is a non-empty-or-empty list of edges.
-- The endpoint types v, w are not tracked.
data Path e = Nil | Cons e (Path e)
\end{lstlisting}

\begin{fidelitynote}
In Lean, \lean{Path G v w} is a \emph{type} indexed by both endpoints; the
type system enforces that consecutive edges match.  In Haskell, this
dependency is erased: a path is just a list of edges, and endpoint
compatibility is a runtime property only.
\end{fidelitynote}


% --- Path.append ---------------------------------------------------
\subsection{Definition: \lean{Path.append}}

\begin{defn}[Path concatenation]
Given paths $\rho : \mathrm{Path}_G(v, w)$ and $\sigma : \mathrm{Path}_G(w, u)$,
their concatenation $\rho \mathbin{+\!\!+} \sigma : \mathrm{Path}_G(v, u)$ is
defined by recursion on $\rho$:
\[
  \mathrm{nil}(v) \mathbin{+\!\!+} \sigma = \sigma, \qquad
  \mathrm{cons}(e, \rho) \mathbin{+\!\!+} \sigma =
    \mathrm{cons}(e,\; \rho \mathbin{+\!\!+} \sigma).
\]
\end{defn}

\paragraph{Lean.}
\begin{lstlisting}[style=lean]
def append : {v w u : G.V} -> Path G v w -> Path G w u -> Path G v u
  | _, _, _, .nil _,    q => q
  | _, _, _, .cons e p, q => .cons e (p.append q)
\end{lstlisting}

\paragraph{Haskell.}
\begin{lstlisting}[style=haskell]
append :: Path e -> Path e -> Path e
append Nil         q = q
append (Cons e p)  q = Cons e (append p q)
\end{lstlisting}


% --- nil_append / append_nil / append_assoc -----------------------
\subsection{Theorems: \lean{Path.nil\_append}, \lean{Path.append\_nil}, \lean{Path.append\_assoc}}

\begin{thm}[Path monoid laws]
For all paths $\rho : \mathrm{Path}_G(v, w)$,
$\sigma : \mathrm{Path}_G(w, u)$, $\tau : \mathrm{Path}_G(u, x)$:
\begin{enumerate}
  \item \textbf{Left unit}: $\mathrm{nil}(v) \mathbin{+\!\!+} \rho = \rho$ (definitional).
  \item \textbf{Right unit}: $\rho \mathbin{+\!\!+} \mathrm{nil}(w) = \rho$.
  \item \textbf{Associativity}: $(\rho \mathbin{+\!\!+} \sigma) \mathbin{+\!\!+} \tau = \rho \mathbin{+\!\!+} (\sigma \mathbin{+\!\!+} \tau)$.
\end{enumerate}
\end{thm}

\paragraph{Lean.}
\begin{lstlisting}[style=lean]
theorem nil_append {v w : G.V} (p : Path G v w) :
    (Path.nil v).append p = p

theorem append_nil : {v w : G.V} -> (p : Path G v w) ->
    p.append (Path.nil w) = p

theorem append_assoc : {v w u x : G.V} ->
    (p : Path G v w) -> (q : Path G w u) -> (r : Path G u x) ->
    (p.append q).append r = p.append (q.append r)
\end{lstlisting}

\paragraph{Haskell.}
\begin{lstlisting}[style=haskell]
-- Left unit  (by definition):  append Nil q = q
-- Right unit (by induction):   append p Nil = p
-- Associativity (by induction):
--   append (append p q) r = append p (append q r)
\end{lstlisting}


%===================================================================
\section{The network directed container}
%===================================================================

% --- Digraph.network -----------------------------------------------
\subsection{Definition: \lean{Digraph.network}}

\begin{defn}[Network container]
For a digraph $G$, the \textbf{network directed container} $\mathcal{N}(G)$ is
the directed container with:
\begin{align*}
  \mathrm{Shape}   &= G.V, \\
  \mathrm{Pos}(v)  &= \textstyle\sum_{w : G.V} \mathrm{Path}_G(v, w), \\
  \mathrm{root}(v) &= (v,\, \mathrm{nil}(v)), \\
  \mathrm{sub}(v, (w, \rho)) &= w, \\
  \mathrm{shift}(v, (w, \rho), (u, \sigma)) &= (u,\, \rho \mathbin{+\!\!+} \sigma).
\end{align*}
The directed-container laws D1, D2, D4 hold definitionally (\lean{rfl});
D3 is the right-unit law and D5 is associativity of \lean{append}.
\end{defn}

\paragraph{Lean.}
\begin{lstlisting}[style=lean]
def network (G : Digraph) : DirectedContainer where
  Shape := G.V
  Pos v := (w : G.V) * Path G v w
  root v := <|v, .nil v|>
  sub _ p := p.1
  shift _ p q := <|q.1, p.2.append q.2|>
  d1 _ := rfl
  d4 _ _ _ := rfl
  d2 _ _ := rfl
  d3 _ p := congrArg (Sigma.mk p.1) (Path.append_nil p.2)
  d5 _ p q r := congrArg (Sigma.mk r.1)
                  (Path.append_assoc p.2 q.2 r.2)
\end{lstlisting}

\paragraph{Haskell.}
\begin{lstlisting}[style=haskell]
-- A directed container (simplified: positions are just paths, not
-- endpoint-paired, since Haskell lacks dependent pairs in the same way).
data Network v e = Network
  { netRoot  :: v -> (v, Path e)
  , netSub   :: v -> (v, Path e) -> v
  , netShift :: v -> (v, Path e) -> (v, Path e) -> (v, Path e)
  }

network :: Digraph v e -> Network v e
network g = Network
  { netRoot  = \v -> (v, Nil)
  , netSub   = \_ p -> fst p
  , netShift = \_ p q -> (fst q, append (snd p) (snd q))
  }
\end{lstlisting}

\begin{fidelitynote}
The five directed-container laws (D1--D5) encode the monoidal structure of path
concatenation.  D1 says the empty path ends at its own start node; D4 says
the endpoint of a concatenated path is the endpoint of the second segment;
D3 is right-unit; D5 is associativity.  D2 is the left-unit law, but must be
stated via a transport along D1 in the general directed-container formalism;
here it reduces to \lean{rfl} because D1 is definitional.
\end{fidelitynote}


%===================================================================
\section{The Sawasdee Cellars wine network}
%===================================================================

% --- Wine.Facility / Wine.Leg / Wine.graph -------------------------
\subsection{Types and graph: \lean{Wine.Facility}, \lean{Wine.Leg}, \lean{Wine.graph}}

\begin{defn}[Sawasdee Cellars network]
The \textbf{Sawasdee Cellars} shipping network (\S2 of the paper) has:
\begin{itemize}
  \item 11 facilities: three vineyards, the Sawasdee export agency, Melbourne
    Port, Thai Port, two distributors, three retail shops;
  \item 10 legs: three vineyard-to-Sawasdee legs (\lean{v1s}, \lean{v2s},
    \lean{v3s}), one Sawasdee-to-port (\lean{smp}), the sea-freight leg
    (\lean{sea}), two port-to-distributor legs (\lean{td1}, \lean{td2}), and
    three distributor-to-retail legs (\lean{d1r1}, \lean{d1r2}, \lean{d2r3}).
\end{itemize}
This is a polytree (hence thin as a directed container).
\end{defn}

\paragraph{Lean.}
\begin{lstlisting}[style=lean]
inductive Facility
  | vineyard1 | vineyard2 | vineyard3
  | sawasdee | melbournePort | thaiPort
  | distributor1 | distributor2
  | retail1 | retail2 | retail3
  deriving DecidableEq, Repr

inductive Leg
  | v1s | v2s | v3s | smp | sea
  | td1 | td2 | d1r1 | d1r2 | d2r3
  deriving DecidableEq, Repr

@[reducible] def graph : Digraph where
  V := Facility
  E := Leg
  src := Leg.src
  tgt := Leg.tgt
\end{lstlisting}

\paragraph{Haskell.}
\begin{lstlisting}[style=haskell]
data Facility = Vineyard1 | Vineyard2 | Vineyard3
              | Sawasdee  | MelbournePort | ThaiPort
              | Distributor1 | Distributor2
              | Retail1 | Retail2 | Retail3

data Leg = V1s | V2s | V3s | Smp | Sea
         | Td1 | Td2 | D1r1 | D1r2 | D2r3
\end{lstlisting}


% --- Wine.network --------------------------------------------------
\subsection{Definition: \lean{Wine.network}}

\begin{defn}[Wine shipping network as a directed container]
\lean{Wine.network} $:= \mathcal{N}(\text{graph})$ is the directed container
obtained by applying Construction~\texttt{con:network} to the Sawasdee Cellars
graph.  Its shapes are the 11 facilities, and its positions at a facility $v$
are all directed routes out of $v$ (including the trivial route).
\end{defn}

\paragraph{Lean.}
\begin{lstlisting}[style=lean]
def network : DirectedContainer := Digraph.network graph
\end{lstlisting}

\paragraph{Haskell.}
\begin{lstlisting}[style=haskell]
wineNetwork :: Network Facility Leg
wineNetwork = network wineGraph
\end{lstlisting}


% --- Sanity checks ------------------------------------------------
\subsection{Sanity checks (examples checked by \lean{rfl})}

\begin{prop}[Route count at Vineyard$_1$]
The position set $\mathrm{Pos}(\mathrm{vineyard1})$ has exactly 9 elements:
the empty route, plus one route to each of the eight facilities reachable from
Vineyard$_1$ (Sawasdee, Melbourne Port, Thai Port, Distributor$_1$,
Distributor$_2$, Retail$_1$, Retail$_2$, Retail$_3$).
\end{prop}

\begin{prop}[Retail shop is a sink]
The position set $\mathrm{Pos}(\mathrm{retail1})$ has exactly 1 element: the
empty route.  Retail shops are sinks in the shipping graph.
\end{prop}

\begin{prop}[\lean{sub} is the route endpoint]
For the full Vineyard$_1 \to$ Retail$_1$ route \lean{fullRoute}:
$\mathcal{N}(\text{graph}).\mathrm{sub}(\mathrm{vineyard1}, \mathrm{fullRoute})
= \mathrm{retail1}$.
\end{prop}

\paragraph{Lean.}
\begin{lstlisting}[style=lean]
example : rVineyard1.length = 9 := rfl
example : rRetail1.length = 1 := rfl
example : network.sub vineyard1 fullRoute = retail1 := rfl
example : network.shift vineyard1 (network.root vineyard1) fullRoute
          = fullRoute := rfl
\end{lstlisting}

\paragraph{Haskell.}
\begin{lstlisting}[style=haskell]
-- These are definitional equalities in Lean (rfl checks),
-- corresponding to properties of the concrete enumerations.
--
-- routeCount Vineyard1 = 9
-- routeCount Retail1   = 1
-- sub Vineyard1 fullRoute = Retail1
-- shift Vineyard1 (root Vineyard1) fullRoute = fullRoute
\end{lstlisting}


%===================================================================
\section{Translation conventions}
%===================================================================

\begin{enumerate}
\item \textbf{Dependent endpoint indices.}  In Lean, \lean{Path G v w} and
  positions $\Sigma_{w} \mathrm{Path}_G(v, w)$ are genuinely dependent types.
  In Haskell, we drop the endpoint tracking and use plain lists or pairs.

\item \textbf{Definitional vs.\ propositional equality.}  Several directed-
  container laws hold definitionally (\lean{rfl}) because the definitions are
  chosen so that the relevant type-level equalities reduce by computation.
  This precision is invisible in the Haskell reading.

\item \textbf{Proof bodies omitted.}  Only definitions and theorem
  \emph{statements} are translated.  Lean checks the proofs; the human
  verifier checks that the \emph{statements} match the intended mathematics.
\end{enumerate}

\end{document}
